\chapter{Implementazione}
\label{chap:implementazione}
%
Conclusa la fase di progettazione, questo capitolo descrive come le scelte architetturali discusse nel capitolo precedente sono state effettivamente realizzate, seguendone la medesima scomposizione. Si parte dallo stack tecnologico adottato e dall'organizzazione del codice (\cref{sec:stack-tecnologico-organizzazione-codice}), per poi approfondire le decisioni implementative relative ai componenti del nucleo del gateway (\cref{sec:principali-decisioni-implementative}) e infine scendere al dettaglio di ciascuna connessione protocollare (\cref{subsec:implementazione-delle-connessioni}).
%
\section{Stack tecnologico e organizzazione del codice}
\label{sec:stack-tecnologico-organizzazione-codice}
%
Definita l'architettura del gateway nei suoi elementi essenziali, questa sezione ne descrive la realizzazione concreta: il linguaggio e gli strumenti adottati, le librerie su cui si appoggiano le singole connessioni, e il modo in cui il codice è organizzato.

\paragraph{Linguaggio e gestione delle dipendenze}
Il gateway è stato sviluppato in Python 3.13, scelto per il supporto nativo ad \texttt{asyncio}. Trattandosi di un componente la cui natura è quasi interamente di I/O -- richieste \ac{HTTP}, sessioni WebSocket, datagrammi \ac{CoAP} e le due direzioni della connessione \ac{MQTT} -- un modello a singolo processo e singolo event loop, in cui ogni operazione di rete è una coroutine concorrente, è risultato ideale a soddisfare il requisito di responsività sotto carico senza dover introdurre concorrenza basata su thread o processi multipli, i quali avrebbero aumentato notevolmente la complessità della soluzione. La gestione delle dipendenze e del processo di build è affidata a \texttt{uv}~\cite{uv}, scelto principalmente per la sua velocità.

\paragraph{Librerie principali}
Attorno a Python gravitano poche librerie, ciascuna responsabile di un solo protocollo o di una sola responsabilità trasversale e capace di fornire un servizio asincrono. FastAPI~\cite{fastapi} è stata scelta per le connessioni \ac{HTTP} e WebSocket poiché le supporta entrambe all'interno della stessa applicazione \ac{ASGI}, eseguita tramite il server Uvicorn~\cite{uvicorn}, generando automaticamente la documentazione \ac{OpenAPI} prevista in fase di progettazione direttamente dalle firme dei router e dai modelli Pydantic~\cite{pydantic}, e offrendo un meccanismo di dependency injection -- riutilizzato per il controllo di salute descritto più avanti -- senza richiedere codice aggiuntivo. Per \ac{CoAP} è stata adottata \texttt{aiocoap}~\cite{aiocoap}, l'implementazione asincrona più matura disponibile per Python, capace di fornire supporto nativo alla modalità Observe. La connessione \ac{MQTT} si appoggia su \texttt{aiomqtt}~\cite{aiomqtt} libreria asincrona che espone un interfaccia tramite \texttt{async}/\texttt{await} invece che tramite callback. Pydantic, infine, attraversa trasversalmente tutte le connessioni: ogni entità del modello di dominio definito in fase di analisi è un modello Pydantic, usato sia per validare e serializzare i payload \ac{MQTT} sia per i corpi delle richieste \ac{HTTP}, i frame WebSocket e i payload \ac{CoAP}, costituendo un'unica fonte di verità per lo schema dei dati indipendentemente dal protocollo con cui transitano.

\paragraph{Containerizzazione e configurazione}
Il gateway viene distribuito tramite un'immagine Docker costruita in due fasi: una fase di build, in cui \texttt{uv} installa le dipendenze bloccate dal lock file e l'interprete Python pinnato, e una fase finale, basata su un'immagine Debian minimale ed eseguita da un utente non privilegiato, che si limita a copiare l'ambiente virtuale già pronto. Il Dockerfile documenta le porte associate ai protocolli di rete coinvolti, in particolare quella su cui è presente il server \ac{ASGI} che ospita le connessioni \ac{HTTP} e WebSocket e quella \ac{UDP} per il server \ac{CoAP}. La configurazione -- l'indirizzo del broker \ac{MQTT} e, opzionalmente, l'indirizzo su cui deve essere avviato il server \ac{CoAP} -- è esternalizzata tramite variabili d'ambiente o file \texttt{.env}, letta da una classe di impostazioni basata su Pydantic; quest'ultima opzione si è rivelata necessaria perché \texttt{aiocoap} non è in grado di effettuare il bind predefinito su macOS, e va quindi sovrascritta manualmente in quell'ambiente.

\paragraph{Organizzazione dei moduli}
Prima di descrivere nel dettaglio le singole decisioni implementative, è utile delineare brevemente il contesto generale: il pacchetto principale si articola in tre rami, oltre al file \texttt{main}.
Questi rami sono: \texttt{/core}: contiene le interfacce e le classi principali, derivate direttamente dalle considerazioni del \cref{sec:architetturale-di-dettaglio-applicativo}, tra cui Connection, Bus, Cache, Container e Monitor, quest'ultimo descritto più dettagliatamente in seguito; \texttt{/connections}: raccoglie le implementazioni di tutte le connessioni; \texttt{/models/topics}: raccoglie, tramite classi Pydantic, tutti i modelli e i relativi payload delle entità di dominio.

%
\section{Principali decisioni implementative}
\label{sec:principali-decisioni-implementative}
%
Definita l'organizzazione dei moduli, è possibile addentrarsi nelle scelte implementative relative ai componenti raccolti nel pacchetto \texttt{/core}, ciascuno dei quali dà forma concreta alle astrazioni discusse nel capitolo di progettazione. I primi tre paragrafi -- Connection, Bus e Cache -- riprendono nell'ordine gli elementi fondamentali dell'architettura interna del gateway già introdotti nella \cref{sec:architetturale-di-dettaglio-applicativo}; il quarto, Container, descrive come queste tre astrazioni vengono concretamente istanziate e collegate tra loro al momento dell'avvio dell'applicazione; il quinto e ultimo, Monitor, introduce infine un componente trasversale non previsto in fase di progettazione, nato dall'esigenza di esporre in modo centralizzato lo stato di salute del gateway. Esauriti questi aspetti comuni a tutto il gateway, la trattazione scende infine al livello delle singole connessioni (\cref{subsec:implementazione-delle-connessioni}), ricalcando lo stesso passaggio dal generale al particolare già adottato in \cref{subsec:progettazione-delle-connessioni} per descrivere, protocollo per protocollo, le scelte progettuali corrispondenti.

\paragraph{Connection}
La classe astratta \texttt{Connection}, definita in \texttt{core/connection.py}, costituisce il punto di estensione del gateway descritto nella \cref{sec:architetturale-di-dettaglio-applicativo}: ogni sotto-modulo che implementa un protocollo, sia esso rivolto verso l'esterno o verso l'interno, eredita da questa classe e ne implementa i due metodi astratti \texttt{start()} e \texttt{stop()}, che ne definiscono il ciclo di vita. Il costruttore riceve tramite dependency injection i riferimenti a \textit{Inbound Bus} e \textit{Outbound Bus}, rendendoli disponibili alle sottoclassi senza che queste debbano occuparsi della loro creazione. Questa interfaccia minimale è sufficiente a garantire l'uniformità richiesta. Ogni connessione, indipendentemente dal protocollo che implementa, può essere avviata e arrestata in modo identico dal resto del sistema. Lo stesso contratto è rispettato anche da \texttt{ExternalConnections} e \texttt{InternalConnections}, due classi composite che a loro volta implementano \texttt{Connection} ma delegano l'effettiva implementazione di \texttt{start()} e \texttt{stop()} all'insieme delle connessioni che aggregano, invocandole concorrentemente tramite \texttt{asyncio.gather}. Questa scelta consente al punto di ingresso dell'applicazione, descritto più avanti in \texttt{Container}, di rimanere del tutto ignaro del numero e della natura delle connessioni effettivamente attive: esso avvia e arresta due soli gruppi, indipendentemente da quanti protocolli essi contengano al loro interno.

\paragraph{Bus}
Coerentemente con quanto discusso nella \cref{sec:architetturale-di-dettaglio-applicativo}, il bus è realizzato attraverso una classe astratta \texttt{Bus}, che espone i due soli metodi \texttt{publish()} e \texttt{subscribe()}, e un'unica implementazione concreta, \texttt{MessageBus}, riutilizzata per entrambe le istanze -- \textit{Inbound} e \textit{Outbound} -- create una sola volta in \texttt{Container}. \texttt{MessageBus} realizza una semantica fire-and-forget: ogni chiamata a \texttt{publish()} itera sui subscriber registrati e, per ciascuno di essi, schedula una coroutine indipendente tramite \texttt{asyncio.create\_task}, restituendo immediatamente il controllo al chiamante senza attendere che la notifica sia stata effettivamente processata. Ogni invocazione di un subscriber è inoltre incapsulata in un blocco \texttt{try/except} che intercetta e registra qualsiasi eccezione, impedendo che il fallimento di un singolo subscriber comprometta gli altri o il publisher stesso. Questa combinazione di concorrenza e isolamento dei fallimenti è ciò che permette al bus di rimanere non bloccante anche quando uno o più subscriber risultano lenti o temporaneamente non disponibili, a costo di rinunciare a garanzie di ordine o di consegna --- un compromesso già discusso e ritenuto accettabile in fase di progettazione.

\paragraph{Cache}
La classe astratta \texttt{Cache} incapsula la logica di sottoscrizione descritta nella \cref{sec:architetturale-di-dettaglio-applicativo}: il proprio costruttore riceve l'\textit{Outbound Bus} e vi si registra immediatamente come subscriber, cosicché ogni sua sottoclasse erediti automaticamente il comportamento di aggiornamento senza doverlo reimplementare. Le sottoclassi devono limitarsi a fornire i due metodi \texttt{get()} e \texttt{\_set()}, rispettivamente per la lettura e per la scrittura nel proprio backing store, oltre ai metodi di ciclo di vita \texttt{start()} e \texttt{stop()}, predisposti per backend che richiedano un'inizializzazione esplicita -- ad esempio l'apertura di una connessione di rete verso un database esterno -- ma lasciati come no-op nella base. L'unica implementazione attualmente presente, \texttt{InMemoryCache}, si appoggia a un semplice dizionario Python indicizzato per topic. La scelta di mantenere l'interfaccia \texttt{Cache} separata dalla sua implementazione, per quanto in questo caso possa apparire eccessivo rispetto agli obbiettivi del progetto, riflette lo stesso principio di disaccoppiamento adottato per le connessioni. Sostituire \texttt{InMemoryCache} con un'implementazione distribuita, ad esempio basata su Redis, per supportare repliche multiple del gateway, richiederebbe di intervenire su un'unica riga in \texttt{Container}, senza alcun impatto sulle connessioni che la utilizzano.

\paragraph{Container}
Tutte le classi concrete del gateway vengono istanziate e collegate fra loro in un unico punto, la classe \texttt{Container}, che funge da composition root dell'applicazione. Nel proprio costruttore, \texttt{Container} crea le due istanze di \texttt{MessageBus}, la \texttt{InMemoryCache}, l'\texttt{HealthMonitor} (descritto nel paragrafo successivo) e il router FastAPI condiviso, per poi passarli, tramite constructor injection, ai due gruppi \texttt{InternalConnections} ed \texttt{ExternalConnections}. Nessun componente del gateway accede a uno stato condiviso globale o importa un singleton al di fuori di questo grafo di dipendenze: ogni classe dichiara esplicitamente, nella propria firma, tutto ciò di cui ha bisogno, rendendo la struttura delle dipendenze leggibile direttamente dal codice. La funzione \texttt{create\_app()}, definita in \texttt{main.py}, riceve un'istanza già completamente configurata di \texttt{Container} e restituisce l'applicazione FastAPI risultante, senza possedere alcuna conoscenza di come i componenti siano stati effettivamente costruiti. Questa separazione tra costruzione e utilizzo è ciò che rende agevole il testing: la suite di test costruisce un proprio \texttt{Container}, sostituendo semplicemente \texttt{MQTTConnection} con un \textit{mock} che simula il comportamento del broker, senza dover replicare o alterare il resto del cablaggio applicativo.

\paragraph{Monitor}
L'ultimo componente del nucleo, \texttt{HealthMonitor}, nasce dall'esigenza di esporre lo stato di salute del gateway senza che questa responsabilità debba essere gestita autonomamente da ciascuna connessione. Ogni connessione che desidera essere monitorata si registra presso l'\texttt{HealthMonitor} indicando se il proprio corretto funzionamento è da considerarsi critico per il sistema nel suo complesso; in cambio riceve un \texttt{HealthHandle}, un oggetto a sola scrittura che espone esclusivamente il metodo \texttt{report()}. Questa asimmetria è intenzionale: la connessione può comunicare il proprio stato, \texttt{UP} o \texttt{DOWN}, ma non può leggere né influenzare quello delle altre, né necessita di sapere se essa stessa sia stata registrata come critica. Il monitor aggrega questi stati in due modi: \texttt{is\_healthy()} restituisce un valore booleano complessivo, vero solo se tutte le connessioni critiche risultano attualmente \texttt{UP}, mentre \texttt{snapshot()} restituisce una struttura dettagliata serializzabile in JSON con lo stato di ciascuna connessione registrata. Nel gateway, la connessione \ac{MQTT} è registrata come critica, poiché un suo malfunzionamento isola completamente il gateway dal dominio interno, mentre la connessione \ac{CoAP} è registrata come non critica. Lo stato viene infine esposto alla chiamata agli endpoint \ac{HTTP}, sfruttando il meccanismo di dependency injection di FastAPI per restituire un errore \texttt{503} qualora una connessione critica risultasse non disponibile.
%
\subsection{Implementazione delle Connessioni}
\label{subsec:implementazione-delle-connessioni}
%
Questa parte descrive come le connessioni progettate in \cref{subsec:progettazione-delle-connessioni} sono state concretamente implementate. Prima di addentrarsi nel dettaglio delle singole connessioni, viene illustrato lo schema generale che ne accomuna l'organizzazione del codice, comune a tutti i protocolli; seguono poi, in altrettanti paragrafi dedicati, le scelte implementative specifiche di \ac{MQTT}, \ac{HTTP}, \ac{CoAP} e WebSocket.

\paragraph{Schema generale}
Il pacchetto \texttt{/connections} rispecchia direttamente la distinzione discussa in \cref{sec:architetturale-rispetto-al-sistema} tra connessioni interne ed esterne: si articola in due sotto-pacchetti, \texttt{externals} e \texttt{internals}, ciascuno contenente la propria classe composita (\texttt{ExternalConnections} e \texttt{InternalConnections}, già descritte nella sezione precedente) e una cartella \texttt{connections/} in cui risiede l'implementazione di ciascun protocollo. All'interno di ciascuna di queste cartelle, ogni protocollo occupa un proprio sotto-pacchetto autonomo (\texttt{mqtt/}, \texttt{http/}, \texttt{coap/}, \texttt{websocket/}), composto da un file \texttt{connection.py} -- che contiene la sottoclasse di \texttt{Connection} e ne rappresenta il punto di ingresso -- ed eventuali cartelle \texttt{routers/}, \texttt{models/} o \texttt{resources/} proprie del protocollo. Questa simmetria nella struttura delle cartelle rende immediatamente riconoscibile, a chiunque esplori il codice, dove cercare l'implementazione di un dato protocollo, e mantiene indipendenti i sotto-pacchetti.

Le due connessioni costruite su FastAPI, ossia \ac{HTTP} e WebSocket, condividono inoltre un meccanismo di composizione dei router: ciascuna riceve, tramite dependency injection, un riferimento al medesimo \texttt{APIRouter} -- creato una sola volta in \texttt{Container} e definito come \texttt{root\_router} -- sul quale registra i propri endpoint tramite \texttt{include\_router()}. In questo modo, l'applicazione FastAPI, istanziata in \texttt{main.py}, non deve conoscere quali protocolli siano effettivamente attivi: si limita a montare un unico router già popolato. La connessione \ac{CoAP}, non basata su \ac{ASGI}, segue invece un meccanismo di composizione analogo ma indipendente, costruendo un oggetto \texttt{Site} di \texttt{aiocoap} popolato con le risorse di propria competenza.

Trasversalmente a tutte le connessioni, infine, i modelli Pydantic definiti in \texttt{models/topic/payloads.py} -- già introdotti nella \cref{sec:stack-tecnologico-organizzazione-codice} -- costituiscono l'unica rappresentazione delle entità di dominio utilizzata: la stessa classe \\\texttt{RobotPosition}, ad esempio, valida il payload \ac{MQTT} in arrivo dal broker, il corpo di una richiesta \texttt{PUT} \ac{HTTP}, il payload di un messaggio WebSocket e quello di una richiesta \ac{CoAP}. Questo evita qualunque duplicazione di schema tra i protocolli e garantisce che un'eventuale modifica a un'entità si propaghi automaticamente a tutte le connessioni che la utilizzano.

\paragraph{MQTT}
La connessione \ac{MQTT}, \texttt{MQTTConnection}, è realizzata attorno a due coroutine concorrenti e indipendenti, \texttt{\_listener} e \texttt{\_publisher}, avviate come task separati al momento dello \texttt{start()}. Il \texttt{\_listener} apre una sessione verso il broker tramite \texttt{aiomqtt.Client}, si sottoscrive al wildcard \texttt{/\#} per ricevere ogni messaggio pubblicato, e per ciascun messaggio in arrivo invoca \texttt{\_dispatch}, che a sua volta tenta di risolvere il topic verso la corrispondente classe Pydantic attraverso \texttt{resolve\_payload\_class} -- la funzione definita in \texttt{models/topic/mappings.py} che confronta il topic contro un insieme di pattern regolari -- e di validare il payload. Un topic non riconosciuto o un payload che fallisce la validazione vengono scartati con un messaggio di log, senza interrompere il flusso degli altri messaggi; in caso di successo, il payload tipizzato viene pubblicato sull'\textit{Outbound Bus}. Il \texttt{\_publisher}, simmetricamente, si sottoscrive all'\textit{Inbound Bus} tramite il metodo \texttt{\_update} e si limita ad accodare ogni coppia topic/payload ricevuta in una \texttt{asyncio.Queue} interna; un ciclo separato consuma questa coda e pubblica ciascun messaggio sul broker tramite lo stesso client \ac{MQTT}. La presenza della coda disaccoppia la velocità con cui l'\textit{Inbound Bus} consegna i messaggi da quella con cui il broker riesce effettivamente ad accettarli.

Entrambi i task condividono la stessa strategia di resilienza: un qualunque errore di connessione -- broker irraggiungibile, connessione interrotta -- viene intercettato, riportato come stato \texttt{DOWN} tramite l'\texttt{HealthHandle} ricevuto in costruzione, e seguito da un nuovo tentativo dopo un'attesa che cresce esponenzialmente da \texttt{BASE\_RECONNECT\_DELAY} (2 secondi) fino a un tetto massimo di \\\texttt{MAX\_RECONNECT\_DELAY} (60 secondi), con l'aggiunta di un piccolo ritardo casuale (\textit{jitter}) per evitare che più istanze del gateway -- nel caso fossero necessarie -- tentino la riconnessione esattamente nello stesso istante.

\paragraph{HTTP}
La connessione \ac{HTTP}, \texttt{HTTPConnection}, non possiede stato né task in background: la sua interezza si esaurisce nel costruttore, dove vengono istanziati e registrati sul \texttt{root\_router} due router, \texttt{RobotRouter} e \texttt{ComputingRouter}, corrispondenti rispettivamente alle entità legate ai singoli robot e a quelle di input dello sciame descritte in \cref{sec:modellazione-delle-entita}. I metodi \texttt{start()} e \texttt{stop()} risultano pertanto no-op: il ciclo di vita effettivo degli endpoint è già governato da FastAPI una volta che l'applicazione viene avviata. Ogni router segue lo stesso schema per ciascuna entità che espone: un endpoint \texttt{GET} che legge dalla cache e restituisce \texttt{404} se la risorsa non è ancora mai stata pubblicata, e un endpoint \texttt{PUT} che valida il corpo della richiesta -- grazie al binding automatico di FastAPI verso il modello Pydantic dichiarato come tipo del parametro -- e lo pubblica sull'\textit{Inbound Bus}, restituendo un oggetto di risposta che riporta sia il topic di destinazione sia il payload effettivamente inoltrato, utile al chiamante per confermare l'esito dell'operazione.

A ogni endpoint viene iniettata una dipendenza tramite FastAPI: una funzione \texttt{check\_health}, definita internamente a \texttt{HTTPConnection}, che interroga l'istanza ricevuta come parametro \texttt{health\_checkable} -- tipizzata come \texttt{HealthReporter}, l'interfaccia a sola lettura descritta nella sezione precedente -- e solleva un'eccezione \texttt{503} corredata dello snapshot dettagliato qualora una connessione critica risulti non disponibile. In questo modo, un client \ac{HTTP} che tenti di leggere o scrivere una risorsa mentre, ad esempio, la connessione \ac{MQTT} è temporaneamente caduta riceve immediatamente un errore esplicito, piuttosto che un valore di cache potenzialmente non più aggiornato. La documentazione \ac{OpenAPI}, generata automaticamente da FastAPI a partire dalle firme dei router e dai modelli di richiesta e risposta, include anche un esempio di tale risposta di errore.

\paragraph{CoAP}
La connessione \ac{CoAP} è strutturalmente la più articolata tra quelle esterne, poiché \texttt{aiocoap} non offre un livello di routing dichiarativo paragonabile a quello di FastAPI: ogni risorsa va costruita come sottoclasse di \texttt{resource.Resource} (o \texttt{resource.ObservableResource}, per supportare Observe) e registrata esplicitamente in un \texttt{Site} al momento dello \texttt{start()}. Per le risorse di input dello sciame, statiche e in numero fisso, sono state definite due classi dedicate, \texttt{FormationResource} e \texttt{LeaderResource}, entrambe derivate da una base comune \texttt{\_BaseComputingResource} che fattorizza la logica di lettura, scrittura e osservazione condivisa tra le due. Per le risorse legate ai robot, il cui numero non è noto a priori, si è invece optato per un'unica \texttt{RobotsSubsite}, registrata sul prefisso \texttt{/robots} e capace di risolvere dinamicamente, a partire dal path della richiesta, sia l'identificativo del robot sia l'endpoint richiesto (\texttt{position}, \texttt{movement}, \texttt{neighbors} o \texttt{sensing}) tramite la tabella \texttt{\_ROBOT\_ROUTES}; aggiungere un nuovo robot non richiede quindi alcuna modifica al codice, mentre aggiungere un nuovo tipo di endpoint richiede solo un'estensione di quella tabella.

Il supporto a Observe, come descritto in \cref{subsubsec:connessione-coap}, è realizzato facendo sottoscrivere ciascuna risorsa all'\textit{Outbound Bus} fin dal proprio costruttore: ogni volta che transita un messaggio il cui topic corrisponde a quello di competenza della risorsa, questa richiama \texttt{trigger()} su ciascuna \texttt{ServerObservation} attualmente registrata, delegando ad \texttt{aiocoap} l'invio della notifica al client osservatore. La gestione delle sottoscrizioni differisce leggermente tra le due famiglie di risorse: \texttt{\_BaseComputingResource}, occupandosi di un solo topic, mantiene un unico insieme di osservazioni, mentre \texttt{RobotsSubsite}, dovendo distinguere tra topic diversi sullo stesso path radice, mantiene un dizionario che associa ciascun topic al proprio insieme di osservatori. In entrambi i casi, una richiesta \texttt{GET} su una risorsa il cui valore non è ancora presente in cache restituisce una risposta \texttt{2.05 Content} con corpo vuoto anziché \texttt{4.04 Not Found}\footnote{il codice 2.05 è l'equivalente CoAP di HTTP~200, mentre 4.04 corrisponde a HTTP~404}, cosicché l'eventuale registrazione Observe associata rimanga valida in attesa del primo aggiornamento e il client non termini la connessione: molti client \ac{CoAP}, infatti, interrompono l'osservazione non appena ricevono una risposta \texttt{4.04 Not Found}, anche se la richiesta \texttt{GET} era stata effettuata correttamente in modalità Observe.

\paragraph{WebSocket}
L'implementazione della connessione WebSocket è incentrata su un'unica classe, \texttt{WebSocketService}, che funge sia da gestore delle connessioni attive sia da gestore del ciclo di vita delle sottoscrizioni, ruoli espressi nel codice tramite due \texttt{Protocol} distinti, \texttt{WebSocketManager} e \texttt{SubscriptionManager}, di cui \texttt{WebSocketService} è l'unica implementazione concreta. Gli endpoint \ac{REST} per la creazione, lettura ed eliminazione delle sottoscrizioni, descritti in \cref{subsubsec:connessione-websocket}, sono realizzati da \texttt{WebsocketSubscriptionRouter} e delegano interamente alla parte \texttt{SubscriptionManager} del servizio. L'endpoint WebSocket vero e proprio, esposto da \texttt{WebsocketEndpointRouter}, delega invece al metodo \texttt{add()}, che accetta la connessione ed entra in un ciclo di ricezione fino alla disconnessione del client o a un errore di validazione. Ogni frame ricevuto viene interpretato tramite un \texttt{TypeAdapter} Pydantic configurato su un'unione discriminata di due soli tipi di messaggio: un \texttt{SubscriptionMessage}, che lega l'identificativo di una sottoscrizione precedentemente creata alla connessione WebSocket corrente -- popolando il dizionario \texttt{\_subscription\_websocket\_dict} -- e un \texttt{PublishMessage}, che consente al client di scrivere direttamente una nuova entità sull'\textit{Inbound Bus} senza dover aprire una richiesta \texttt{PUT} separata. In direzione opposta, il servizio si sottoscrive all'\textit{Outbound Bus} al momento dello \texttt{start()}: per ogni messaggio in arrivo il cui topic compare tra quelli sottoscritti, viene inviato un \texttt{SubscriptionUpdateMessage} a ciascun WebSocket attualmente legato a una sottoscrizione di quel topic, con gli invii eseguiti concorrentemente e gli eventuali fallimenti di singole connessioni isolati, in modo che la disconnessione di un client non comprometta la consegna agli altri. Infine, ogni connessione WebSocket attiva viene tracciata come task in un insieme dedicato, cancellato in blocco al momento dello \texttt{stop()}, garantendo uno spegnimento ordinato di tutte le sessioni aperte.
%